\documentclass{article}
\setlength\parindent{0pt}
\pagenumbering{gobble}

\usepackage{geometry}
\geometry{papersize={210mm,297mm}}
\geometry{left=9mm,right=9mm,top=9mm,bottom=9mm}

\usepackage{fontspec}
\setmainfont{DejaVu Sans}[Scale=0.9]

\usepackage{xeCJK}
\CJKfamily{zhsong}

\usepackage[svgcolors]{xcolor}
\usepackage{longfbox}
\usepackage{epstopdf}
\usepackage{float}
\usepackage{xpatch}
\usepackage{tipa}
\usepackage[export]{adjustbox}
\usepackage{enumitem}
\usepackage{mdframed}

\usepackage{multirow}
\usepackage{tabularx}
\renewcommand\tabularxcolumn[1]{m{#1}} % for vertical centering text in X column

\usepackage{multicol}
\setlength{\columnsep}{3mm}

\usepackage[explicit,compact]{titlesec}
\titleformat{\section}{\normalfont\bfseries}{}{0pt}{【#1】}

% 啊！万能的 StackExchange！
% https://tex.stackexchange.com/questions/475466/latex-three-column-layout-merging-two-of-them-at-the-begining
%
\newlength{\abstractwidth}
\newlength{\columnshrink}
\newsavebox{\twocolinsert}
%
\makeatletter
\newlength{\resized@col}
\newcounter{column@count}
%
\xpatchcmd{\multi@column@out}{
	\process@cols\mult@gfirstbox{%
		\setbox\count@
		\vsplit\@cclv to\dimen@
		\set@keptmarks
		\setbox\count@
		\vbox to\dimen@
		{\unvbox\count@ \ifshr@nking\vfilmaxdepth\fi}%
	}%
}{
	\process@cols\mult@gfirstbox{%
		\global\advance\c@column@count\@ne
		\resized@col\dimen@%
		\ifnum\c@column@count=\tw@
				\advance\resized@col-\columnshrink
		\fi%
		\setbox\count@
		\vsplit\@cclv to\resized@col
		\set@keptmarks
		\setbox\count@
		\vbox to\dimen@{
			\ifnum
				\c@column@count=\tw@ \vspace*{\columnshrink}
			\fi
			\unvbox\count@
			\ifshr@nking\vfilmaxdepth\fi
		}%
	}%
}{\typeout{Success}}{\typeout{Failure}}
\makeatother

% for designing header
\newsavebox\mysavebox
\newenvironment{imgminipage}[2][]{%
   \def\imgcmd{\includegraphics[width=\wd\mysavebox, height=\dimexpr\ht\mysavebox+\dp\mysavebox\relax, #1]{#2}}%
   \begin{lrbox}{\mysavebox}%
   \begin{minipage}%
}{%
   \end{minipage}
   \end{lrbox}%
   \sbox\mysavebox{\setlength{\fboxrule}{0pt}\fbox{\usebox\mysavebox}}%
   \mbox{\rlap{\raisebox{-\dp\mysavebox}{\imgcmd}}\usebox\mysavebox}%
}


\begin{document}
\begin{multicols*}{3}

	\setlength{\abstractwidth}{2\linewidth}
	\addtolength{\abstractwidth}{\columnsep}
	\savebox{\twocolinsert}{\begin{minipage}{\abstractwidth}
		\noindent 核准日期：2020年10月31日
		\newline 修改日期：2019年07月06日，2020年10月31日
		\newline

		% \begin{mdframed}[leftline=false, rightline=false, innertopmargin=0pt, innerbottommargin=0pt, innerrightmargin=0pt, innerleftmargin=2em]
		% 	\includegraphics[width=0.15\abstractwidth, valign=m]{assets/debian-text.eps}
		% 	\hfill
		% 	\begin{imgminipage}{assets/header-background.eps}[t]{0.7\abstractwidth}
		% 		\Large \textbf{盒装安装媒介说明书}

		% 		\normalsize 请仔细阅读说明书并在管理员指导下使用
		% 	\end{imgminipage}
		% \end{mdframed}

		\includegraphics[width=\abstractwidth]{assets/header.eps}


		\begin{mdframed}[hidealllines=true, innerbottommargin=.5em, innertopmargin=0pt]
			\sffamily
			
			{\centering 警示语 \par}
		
			无论是否与其它操作系统合用，安装 Debian 均存在丢失磁盘上所有内容的风险。（参见【不良反应】）
			
			很多软件是与多媒体相关的，特别是允许回放和提供音频、视频操作或类似功能的软件，不被 Debian 包含，是因为它在世界的某个地方被认为是非法的。本说明不应该被认为是有资格的法律建议的替代；法律建议只能咨询律师来获得。
		
			一些硬件制造商拒绝告诉我们如何给他们的硬件写驱动程序。另一些则要求签署不公开的契约才能接触文档，以阻止我们发布驱动程序源代码，这一自由软件的核心内容。由于我们未被授权使用这些文档，造成它们无法在 Linux 下工作。
		\end{mdframed}
	\end{minipage}}
	\setlength{\columnshrink}{\ht\twocolinsert}
	\addtolength{\columnshrink}{\dp\twocolinsert}
	\noindent\usebox{\twocolinsert}


	\section*{发行版名称}
	\begin{tabularx}{\linewidth}{@{}ll@{}}
		通用名称： & Debian \\
		正式名称： & GNU/Linux Debian \\
		英文音标： & \textipa{["dEbi@n]} \\
	\end{tabularx}


	\section*{内容}

	完全由自由软件组成的类 UNIX 操作系统，其包含的多数软件使用 GNU 通用公共许可协议授权，并由 Debian 计划的参与者组成团队对其进行打包、开发与维护。


	\section*{性质}

	本系统为采用 GNU/Linux内核的操作系统，安装后可由 UEFI 或 Legacy 引导方式启动。

		
	\section*{适应平台}

	支持使用 ARM、IA-32、IA-64、x86-64、PowerPC、ppc64el、ARM、MIPS 等架构的计算机、服务器和嵌入式设备。

	本包装盒中的安装媒介适用于何平台以实际为准。


	\section*{规格}
	
	1 个 安装媒介


	\section*{用法}

	使用 USB 设备引导。

	启动方式根据硬件调整，一般使用 UEFI。
	
	根据硬件性能和个人需要，调整安装方式：一般而言，使用图形化界面进行安装；否则选择基于 ncurses 命令行下的安装。
	
	安装好基本系统并设置完成软件包后，即可酌情选择桌面环境和桌面管理器。
		

	\section*{不良反应}

	Debian 有一个对用户和开发者所提交的软件缺陷报告进行归档管理的缺陷跟踪系统，英文缩写为 BTS。每个软件缺陷报告都被授予一个编号并且被长期跟踪，直到它被标记为已修复。
	
	可以在 https://www.debian.org/Bugs/ 处获得这个文件的拷贝。


	\section*{注意事项}
	\begin{itemize}[leftmargin=*]

		\item 满足最低的硬件要求
		
		基于您的需求，也许可以用低于下面表格所列的配置装上系统。但是，如果无视这些建议的话，多数用户会安装失败。

		{\small\begin{tabularx}{\linewidth}{|X|X|X|X|}
			\hline
			类别 & RAM\newline (最低) & RAM\newline (推荐) & 硬盘 \\
			\hline
			无桌面 & 128MB & 512MB & 2GB \\
			\hline
			有桌面 & 256MB & 1GB & 10GB \\
			\hline
		\end{tabularx}}

		Pentium 4、1GHz 的系统是桌面系统的最低推荐配置。

		\item 需要固件的设备
		
		除了需要设备驱动程序，有些硬件还要在使用之前加载固件（firmware）或微码（microcode）。

	\end{itemize}


	\section*{禁忌}

	在操作过程中出现或即将出现下列任何一种情况，请立即停止操作，并准备好系统恢复。

	\begin{itemize}[leftmargin=*]
		\setlength{\itemsep}{0pt}
		\setlength{\parskip}{0pt}
		\setlength{\parsep}{0pt}

		\item 以根权限在根目录下执行递归删除
		\item 未确认设备名即使用 dd 命令
		\item 未确认设备名即执行格式化
		\item 未确认操作即使用命令重定向
		\item 未经确认即执行来自网络的脚本
		\item 长期在散热不良的设备上高负载使用
	\end{itemize}


	\section*{系统相互作用}
	\begin{itemize}[leftmargin=*]
		\setlength{\parindent}{0pt}

		\item 与 Windows 的相互作用
		
		当您有双引导时，若另一个操作系统与 Windows 访问相同的文件系统，这就有可能会导致问题和数据丢失。在这种情况下，文件系统的真实状态可能与 Windows 8 认为在“启动”之后的情况不同，并且可能在进一步写入文件系统时导致文件系统损坏。因此，在双引导设置中，为了避免文件系统损坏，必须在 Windows 中禁用“快速启动”功能。
		
		在罕见情况中已观察到，在使用 Windows 进行系统更新时，可能会出现重新启动后 Grub 引导被破坏从而导致 Debian 无法启动的情况。同时，若在安装过程中将引导信息写入与 Windows 所在的物理磁盘 MBR 内，将导致后者无法正常启动。

		\item 与其他 GNU/Linux 发行版的相互作用尚不明确。

	\end{itemize}


	\section*{贮藏}

	-40℃\textasciitilde +70℃

	妥善贮藏所有安装媒介，勿使不会安装的人员触及。


	\section*{包装}

	装有 GNU/Linux Debian 安装镜像的、兼容 USB 3.0/2.0 协议的大容量存储设备。
	
	1 枚/盒


	\section*{有效期}

	完整支持：36 个月

	长期支持：60 个月


	\section*{执行标准}
	\begin{tabularx}{\linewidth}{@{}ll@{}}
		\multirow{4}{*}{}{开源许可证：} & GNU GPL（主要）\\
		~ & GNU LGPL \\
		~ & BSD \\
		~ & 以及其他授权条款 \\
	\end{tabularx}


	\section*{批准文号}

	说明书使用 CC-BY-SA 3.0 协议授权。


	\section*{生产单位}

	Debian 计划


	\section*{说明书}
	\begin{tabularx}{\linewidth}{@{}ll@{}}
		\multirow{2}{*}{}{编审：} & @SteveCharlesYang\\
		~ & @moesoha \\
		图形: & @YJBeetle\\
		GitHub: & moesoha/debian-media-box\\
	\end{tabularx}


	\vfill
	\begin{flushright}
		GNU/Linux Debian - 10.0\linebreak 20190706
		\linebreak
		\newline
		\includegraphics[width=13mm]{assets/debian-logo.eps}
	\end{flushright}

\end{multicols*}
\end{document}
